\section{Explicit Calculation at Second Order}

\subsection{Duhamel's formula}

\textcolor{red}{Motivate the study of this system}.

The state space for a system of N identical fermions in a box of side length L with periodic
boundary conditions is the space of completely antisymmetric $L^2$-functions on the torus
$\mathbb{T}^{3N} = (\mathbb{R}/2\pi\mathbb{Z})^{3N}$:

\begin{align*}
	L^2_a(\mathbb{T}^{3N}) := \{\psi\in L^2(\mathbb{T}^{3N})
	: \psi(x_{\sigma(1)}, ..., x_{\sigma(N)})
	= \mathrm{sgn}(\sigma)\psi(x_1, ..., x_N) \,\forall \sigma \in S_N\}
\end{align*}

The Hamiltonian for weakly interacting particles can be written as a sum of
the kinetic energy terms plus a sum of two-body interaction terms:
\begin{align*}
	H = \hbar^2 \sum_{i=1}^{N}(-\Delta_{x_i}) + \lambda \sum_{1 \leq i <
	j \leq N}V(x_i - x_j)
\end{align*}
where $V : \mathbb{R}^3 \to \mathbb{R}$ is a smooth function
with period $L$ in all variables and such that $V(x_i - x_j) = V(|x_i - x_j|)$,
in particular this means that $\hat{V}(-k) = \hat{V}(k)$, and this fact will become useful
in evaluating the second order perturbation to the ground state energy of the system. The
scaling factor $\lambda \in \mathbb{R}$ is assumed to be small. Working in the grand canonical ensemble
requires to lift this operator to the fermionic Fock space, and in second quantization
formalism this means that the Hamiltonian will be written as
\begin{align*}
	H = H_0 + V = \sum_{k \in \mathbb{Z}^3} \hbar^2|k|^2 a^*_k a_k
	+ \sum_{k,p,q \in \mathbb{Z}^3} \lambda\hat{V}(k) a^*_{p-k} a^*_{q+k} a_q a_p
\end{align*}
where $\hat{V}(k)$ are the Fourier coefficients of $V$ and $a^\#_k$ are creation/annihilation
operators on the Fock space that create/annihilate states of definite momentum $k \in \mathbb{Z}^3$.

Since we want to
compute the second order contribution to the ground state energy, we will calculate the
free energy of the system and take its limit as the temperature goes to zero,
that is $\beta \to \infty$.
Recall that at thermal equilibrium the expectation value of any observable is
given by the Gibbs state functional
\begin{align*}
	\langle\cdot\rangle = \frac{\mathrm{Tr}(\cdot \, e^{-\beta H})}
	{\mathrm{Tr}(e^{-\beta H})}.
\end{align*}

Because the free energy is given by the logarithm of the
partition function, we can make use of a useful formula that allows to truncate
the quantity $\mathrm{log}(\mathrm{Tr}(e^{-\beta H}))$ to second order contributions
in the scaling factor $\lambda$.

\begin{proposition}
	Given $H = H_0 + V$ as defined above, the following approximation holds \textcolor{red}{up to terms
	of order $O(...)$ (consider putting $\lambda$ outside of $V$, i.e. $H = H_0 + \lambda V$, to make
	this explicit)}.
	\begin{align}
		\mathrm{log}\,\mathrm{Tr}\,e^{-\beta H}
		= \mathrm{log}\,\mathrm{Tr}\,e^{-\beta H_0} - \beta \langle V \rangle_{H_0}
		\int_0^\beta\dx\tau \int_0^\tau\dx\tau' \left(
			\langle V(\tau')V(\tau) \rangle_{H_0} - \langle V(\tau') \rangle_{H_0}\langle V(\tau) \rangle_{H_0}
		\right)
	\end{align}
\end{proposition}
\begin{proof}
	Start by rewriting the operator $e^{-\beta H}e^{\beta H_0}$ using the
	fundamental theorem of calculus:
	\begin{align*}
		e^{-\beta H}e^{\beta H_0}
		&= \mathbbm{1}
			+ \int_0^\beta\dx\tau
			(\frac{\dx}{\dx t}e^{-t H}e^{t H_0})(\tau)\\
		&= \mathbbm{1}
			+ \int_0^\beta\dx\tau
			e^{-\tau H}(-H+H_0)e^{\tau H_0}\\
		&= \mathbbm{1}
			+ \int_0^\beta\dx\tau
			e^{-\tau H}e^{\tau H_0}e^{-\tau H_0}(-V)e^{\tau H_0}\\
		&= \mathbbm{1}
			- \int_0^\beta\dx\tau
			e^{-\tau H}e^{\tau H_0}V(\tau) = \star
	\end{align*}

	\textcolor{red}{The integral converges because of Stone's theorem on strongly continuous
	unitary groups}.
	Now we can see that in the integral we have the same term that appears in the
	l.h.s. but at a generic temperature $\tau$. This means that we can repeat
	the same steps to push this operator to higher and higher "orders". Iterating this procedure
	just once gives the following:
	\begin{align*}
		\star &= \mathbbm{1}
				- \int_0^\beta\dx\tau
				\left(
				\mathbbm{1} + \int_0^\tau\dx\tau'e^{-\tau' H}e^{\tau' H_0}V(\tau')
				\right)V(\tau) \\
			  &= \mathbbm{1} - \int_0^\beta\dx\tau V(\tau)
				+ \int_0^\beta\dx\tau \int_0^\tau\dx\tau' V(\tau')V(\tau)
				+ \mathrm{R}(V,\beta).
	\end{align*}
	where $\mathrm{R}(V,\beta)$ is the remainder of this expression \textcolor{red}{and is of order $O(...)$}.
	Dropping the remainder term and truncating the expression,
	we can invert this relation to isolate $e^{-\beta H}$ and take the trace. Since the trace
	is a series of matrix elements \textcolor{red}{we can move it inside the integral}
	to get:
	\begin{align*}
		\mathrm{Tr} \, e^{-\beta H}
		&=
			\mathrm{Tr} \, e^{-\beta H_0}
			- \int_0^\beta\dx\tau\, \mathrm{Tr}(V(\tau)e^{-\beta H_0})
			+ \int_0^\beta\dx\tau\int_0^\tau\dx\tau'\, \mathrm{Tr}(V(\tau')V(\tau)e^{-\beta H_0})\\
		&= (\mathrm{Tr}\,e^{-\beta H_0})\left(
			1 - \int_0^\beta\dx\tau\,\langle V(\tau) \rangle_{H_0}
			+ \int_0^\beta\dx\tau\int_0^\tau\dx\tau'\, \langle (V(\tau')V(\tau) \rangle_{H_0}
			\right).
	\end{align*}

	We can now take the logarithm on both sides. Recall the Taylor expansion
	$\mathrm{log}(1+x) \sim x - x^2/2\,$ as $x \to 0$ and apply it to the r.h.s.,
	but note that the first integral gives two contributions (one at first
	order and one at second order) while the second integral only contributes to the second order.
	In the end, we are left with:
	\begin{align}
		\label{eq:duhamel}
		\mathrm{log}\,\mathrm{Tr}\,e^{-\beta H}
		&= \mathrm{log}\,\mathrm{Tr}\,e^{-\beta H_0}
			- \int_0^\beta\dx\tau\,\langle V(\tau) \rangle_{H_0} \nonumber \\
		&\quad + \int_0^\beta\dx\tau\int_0^\tau\dx\tau'\, \langle (V(\tau')V(\tau) \rangle_{H_0}
			-\frac{1}{2} \left( \int_0^\beta\dx\tau \langle V(\tau) \rangle_{H_0} \right)^2 \nonumber \\
		&= \mathrm{log}\,\mathrm{Tr}\,e^{-\beta H_0}
			- \beta \langle V \rangle_{H_0} \nonumber \\
		&\quad+ \int_0^\beta\dx\tau\int_0^\tau\dx\tau'\,\biggl(
				\langle V(\tau')V(\tau) \rangle_{H_0} - \langle V(\tau') \rangle_{H_0} \langle V(\tau) \rangle_{H_0}
			\biggr)
	\end{align}
	where in the last step we performed the first integral and rewrote the square as a double integral,
	then taking into account the symmetry under $\tau \longleftrightarrow \tau'$
	we could remove the factor $1/2$ by integrating only over the lower triangle $\{(\tau,\tau')\in[0,\beta]^2 : \tau' \leq \tau\}$.
\end{proof}

\subsection{Evaluating expectation values}

To evaluate the expectation values, which are expectations of polynomials in field operators,
we can turn to Grassmann variables. Indeed, \textcolor{red}{by introducing a cutoff $\Lambda$}
we see that the possible momentum states form a finite set of M independent modes:
\begin{align*}
	\mathcal{D}_\Lambda := \{...\}.
\end{align*}

Let then $\boldsymbol{\psi} = (\psi_1, ..., \psi_M)$ and
$\boldsymbol{\bar\psi} = (\bar\psi_1, ..., \bar\psi_M)$ be the Grassmann variables
associated to the families of operators $a$ and $a^*$ respectively.
Having introduced the cutoff, the Hamiltonian can now be written as a simple
finite linear combination of creation/annihilation operators:
\begin{align*}
	H_0 = \sum_{k \in \mathcal{D}_\Lambda} \epsilon_k a^*_k a_k
\end{align*}

We begin the calculation by building the gaussian measure for the expectations
$\langle\,\cdot\,\rangle_{H_0}$.
Note that $\langle V \rangle_{H_0}$ contains field operators at the same temperature, while
the expectations inside the double integral contain field operators at different temperatures,
which means the two-point funcitons that appear in the contractions of the respective
diagrams are going to be different.

\begin{lemma}
	\label{two-point-function}
	The two-point function for this system has the following form:
	\begin{equation}
		G_{k,k'}(\tau,\tau')
		= \langle a^*_k(\tau) a_{k'}(\tau') \rangle_{H_0}
		= \delta_{k,k'} n_{\mathrm{F}}(\epsilon_k) e^{-(\tau-\tau')\epsilon_{k}}
	\end{equation}
	where $n_\mathrm{F}(\epsilon) = (e^{\beta\epsilon}+1)^{-1}$ is the Fermi-Dirac distribution.
\end{lemma}
\begin{proof}
	We are going to divide the proof in two parts: first we will find an expression for the
	operators $a^\#(\tau)$, as this will allow us to write the two-point function in terms of
	$\langle a^*_k a_k' \rangle$, then we will compute this last expectation to find the
	Fermi-Dirac distribution. Recall that in the Heisenberg picture with imaginary time one has
	that for an operator $A$ its evolution at temperature $\tau$ is given by
	$A(\tau) = e^{-\tau H_0}Ae^{\tau H_0}$. \textcolor{red}{check domains}. Once again we will
	make use of the fundamental theorem of calculus to write
	\begin{align*}
		e^{-\tau H_0}a_k e^{\tau H_0}
		&= a_k + \int_0^\tau \dx\lambda\, \frac{\dx}{\dx\lambda}\left( e^{-\lambda H_0}a_k e^{\lambda H_0} \right) \\
		&= a_k + \int_0^\tau \dx\lambda\, e^{-\lambda H_0}(-H_0 a_k + a_k H_0)e^{\lambda H_0} \\
		&= a_k + \int_0^\tau \dx\lambda\, e^{-\lambda H_0}[a_k,H_0]e^{\lambda H_0} \\
		&= a_k + \int_0^\tau \dx\lambda\, e^{-\lambda H_0}\sum_{k'\in\mathcal{D}_\Lambda}\epsilon_{k'}[a_k,a^*_{k'} a_{k'}]e^{\lambda H_0}
		= \star.
	\end{align*}

	Recall that, for operators $A$, $B$ and $C$ defined on a common domain, it holds that
	\begin{align*}
		[A,BC] = \{A,B\}C - B\{A,C\}
	\end{align*}
	as can be shown by explicitly writing the expressions for the anticommutators. Using this
	identity we rewrite the commutator that appears in the last expression in terms of the
	canonical anticommutator as follows:
	\begin{align*}
		\star
		&= a_k + \int_0^\tau\dx\lambda\, e^{-\lambda H_0}
			\sum_{k'\in\mathcal{D}_\Lambda}\epsilon_{k'}(\{a_k,a^*_{k'}\}a_{k'} - a^*_{k'}\{a_k,a_{k'}\})
			e^{\lambda H_0} \\
		&= a_k + \int_0^\tau\dx\lambda\, e^{-\lambda H_0}\epsilon_{k}a_{k}e^{\lambda H_0} \\
		&= a_k + \epsilon_k\int_0^\tau\dx\lambda\,a_k(\lambda)
	\end{align*}
	where we used that $\{a_k,a^*_{k'}\}=\delta_{k,k'}$ and that $\{a_k,a_{k'}\}=0$. We now
	have an expression for $a_k(\tau)$ in terms of its integral, and if we take the derivative
	with respect to $\tau$ we get the following Cauchy problem:
	\begin{align*}
		\begin{cases}
			\frac{\dx}{\dx\tau}a_k(\tau) = \epsilon_k a_k(\tau) \\
			a_k(0) = a_k
		\end{cases}.
	\end{align*}
	
	By the \textcolor{red}{theorem of existence and uniqueness of the solution to an ODE on
	a closed interval}, it suffices to find one explicitly. This is a well known differential
	equation and is easy to check that the following is indeed a solution
	\begin{align*}
		a_k(\tau) = e^{\tau\epsilon_k}a_k(0) = e^{\tau\epsilon_k}a_k.
	\end{align*}

	The procedure is completely analogous in the case of $a^*_k(\tau)$ and the result differs
	only by a sign in the exponent:
	\begin{align*}
		a^*_k(\tau) = e^{-\tau\epsilon_k}a^*_k(0) = e^{-\tau\epsilon_k}a^*_k.
	\end{align*}

	By linearity, we can take the scalar exponential factors out of the expectation value to
	write the two-point function in terms of the quantity $\langle a^*_k a_{k'}\rangle$:
	\begin{align*}
		\langle a^*_k(\tau) a_{k'}(\tau') \rangle_{H_0}
		= e^{-\tau\epsilon_k}e^{\tau'\epsilon_{k'}} \langle a^*_k a_{k'} \rangle_{H_0}.
	\end{align*}

	It remains to determine the expectation $\langle a^*_k a_{k'}\rangle$ and find the Fermi-Dirac
	distribution. Recall that the spectrum of the number density operators $a^*_k a_k$ is
	$\mathbb{N}_0$, and in the fermionic case the occupation number basis decomposes the Fock
	space like we discussed in Proposition \ref{occupation-number-basis}. Using the cutoff on high
	momentum, the trace in the Gibbs state functional becomes a sum over a finite amount of
	states, therefore let
	$\{\vert\boldsymbol{n}\rangle = \vert n_1\rangle\otimes\cdots\otimes\vert n_M\rangle
	: n_i\in\{0,1\}\;\forall i=1,\ldots,M\}$
	be the complete orthonormal set of eigenvectors of $M$ independent number density operators
	(in our case, $M=\vert\mathcal{D}_\Lambda\vert$). Calling $Z = \mathrm{Tr}\,e^{-\beta H_0}$
	and using the momenta $k\in\mathcal{D}_\Lambda$ as indices for the basis states labels $n_k$,
	we write
	\begin{align*}
		\langle a^*_k a_{k'} \rangle_{H_0}
		&= \frac{\mathrm{Tr}(a^*_k a_{k'}e^{-\beta H_0})}{\mathrm{Tr}\,e^{-\beta H_0}}
		= \frac{1}{Z} \sum_{\boldsymbol{n}\in\{0,1\}^M}\langle\boldsymbol{n}\vert e^{-\beta H_0}a^*_k a_{k'}\vert\boldsymbol{n}\rangle\\
		&= \delta_{k,k'} \frac{1}{Z} \sum_{\boldsymbol{n}\in\{0,1\}^M}\langle\boldsymbol{n}\vert a^*_k a_k \vert\boldsymbol{n}\rangle\prod_{p\in\mathcal{D}_\Lambda}e^{-\beta\epsilon_p n_p}\\
		&= \delta_{k,k'} \frac{1}{Z} \sum_{\boldsymbol{n}\in\{0,1\}^M} n_k e^{-\beta\epsilon_k n_k} \prod_{p\neq k} e^{-\beta\epsilon_p n_p}\\
		&= \delta_{k,k'} \frac{1}{Z} \left(\sum_{n_k=0}^1 n_k e^{-\beta\epsilon_k n_k}\right)\prod_{p\neq k}\sum_{n_p=0}^1 e^{-\beta\epsilon_p n_p}\\
		&= \delta_{k,k'} \frac{1}{Z} e^{-\beta\epsilon_k}\prod_{p\neq k}(1 + e^{-\beta\epsilon_p})
	\end{align*}
	where we used the fact that $e^{-\beta H_0}$ \textcolor{red}{is bounded} and diagonal
	in the occupation number basis $\{\vert\boldsymbol{n}\rangle\}$, and that
	$a^*_k a_k\vert\boldsymbol{n}\rangle=n_k\vert\boldsymbol{n}\rangle$. Using many of the
	same steps for the denominator, we find that $Z = \prod_{p\in\mathcal{D}_\Lambda}(1 + e^{-\beta\epsilon_p})$,
	thus the expression simplyfies to
	\begin{align*}
		\langle a^*_k a_{k'}\rangle_{H_0}
		&= \delta_{k,k'} \frac{e^{-\beta\epsilon_k}\prod_{p\neq k}(1 + e^{-\beta\epsilon_p})}{\prod_{p}(1 + e^{-\beta\epsilon_p})}\\
		&= \delta_{k,k'} \frac{e^{-\beta\epsilon_k}}{1 + e^{-\beta\epsilon_k}}\\
		&= \delta_{k,k'} \frac{1}{e^{\beta\epsilon_k}+1}.
	\end{align*}
	This is exactly the Fermi-Dirac distribution, and the thesis follows immediately.
\end{proof}

With this quantity we can go over to the computation via Feynman diagrams. Recall the definition
for the expectation of a function of Grassmann variables
\begin{align*}
	\mathbb{E}_G(\,\cdot\,) = \int P_G(\dx\psi)(\,\cdot\,)\, , \quad \text{where}\;
	P_G(\dx\psi) = \frac{1}{\mathrm{det}\,G} \left( \prod_{j=1}^N \dx\psi_j\dx\bar\psi_j \right)
	e^{-\boldsymbol{\bar\psi}G^{-1}\boldsymbol{\psi}}
\end{align*}
and the graphical representation of the interaction term via a vertex:
\begin{align*}
	V = \lambda \sum_{k,p,q} \hat{V}(k) \bar\psi_{p-k} \bar\psi_{q+k} \psi_q \psi_p = \;
	\begin{tikzpicture}[scale=.8, baseline=(current bounding box.center)]
		\node (a) at (-1, 1) {} node[scale=0.7] at (-1.1, 0.5) {$p-k$};
		\node (b) at ( 1, 1) {} node[scale=0.7] at ( 1.1, 0.5) {$q+k$};
		\node (c) at (-1,-1) {} node[scale=0.7] at (-0.9,-0.5) {$q$};
		\node (d) at ( 1,-1) {} node[scale=0.7] at ( 0.9,-0.5) {$p$};
		\begin{scope}[very thick, decoration={
				markings,
				mark=at position 0.6 with {\arrow{>}}}
			]
			\draw[black, thick, postaction={decorate}] (0,0) -- (a);
			\draw[black, thick, postaction={decorate}] (0,0) -- (b);
		\end{scope}
		\begin{scope}[very thick, decoration={
				markings,
				mark=at position 0.6 with {\arrow{<}}}
			]
			\draw[black, thick, postaction={decorate}] (0,0) -- (c);
			\draw[black, thick, postaction={decorate}] (0,0) -- (d);
		\end{scope}
		\filldraw[black] (0,0) circle (2pt); %node[anchor=west]{Intersection point};
	\end{tikzpicture}.
\end{align*}

With these we can evaluate the expectations that show up in Duhamel's formula separately. Let's start
with $\langle V \rangle_{H_0}$. This is the expectation of a sum with terms $\bar\psi\bar\psi\psi\psi$,
exactly the interaction term, and so the result will be a sum of all possible (connected) diagrams
with one vertex. Omitting the temperatures when writing two-point functions and calling
$n_k = n_\text{F}(\epsilon_k)$ the calculation is as follows:
\begin{align*}
	\langle V \rangle_{H_0}
	&= \mathbb{E}_G \biggl(
		\lambda \sum_{k,p,q} \hat{V}(k) \bar\psi_{p-k} \bar\psi_{q+k} \psi_q \psi_p
		\biggr)
	= \sum_{\Gamma \in \mathcal{G}_{\text{C},1}} \text{val}(\Gamma) \\
	&= \textcolor{red}{draw \ feynman \ diagrams} \\
	% \begin{tikzpicture}[scale=.7, baseline=(current bounding box.center)]
	% 	\node[shape=coordinate] (a) at (-1, 1) {};
	% 	\node[shape=coordinate] (b) at ( 1, 1) {};
	% 	\node[shape=coordinate] (c) at (-1,-1) {};
	% 	\node[shape=coordinate] (d) at ( 1,-1) {};
	% 	\begin{scope}[very thick, decoration={
	% 			markings,
	% 			mark=at position 0.5 with {\arrow{>}}}
	% 		]
	% 		\draw[black, thick] (0,0) -- (a);
	% 		\draw[black, thick] (0,0) -- (b);
	% 	\end{scope}
	% 	\begin{scope}[very thick, decoration={
	% 			markings,
	% 			mark=at position 0.5 with {\arrow{<}}}
	% 		]
	% 		\draw[black, thick] (0,0) -- (c);
	% 		\draw[black, thick] (0,0) -- (d);
	% 	\end{scope}
	% 	\filldraw[black] (0,0) circle (2pt); %node[anchor=west]{Intersection point};
	% 	% \draw[black, thick, postaction={decorate}] (a) .. controls (-2,-.5) and (-.5,-2) .. (d);
	% \end{tikzpicture}.
	&= \sum_{k,p,q} \hat{V}(k) G_{p-k,p}G_{q+k,q} - \sum_{k,p,q} \hat{V}(k) G_{p-k,q}G_{q+k,p} \\
	&= \sum_{k,p,q} \bigl(\hat{V}(k) \delta_{k,0} n_p n_q - \hat{V}(k) \delta_{k,p-q} n_p n_q \bigr)\\
	&= \sum_{p,q} (\hat{V}(0) - \hat{V}(p-q)) n_p n_q
\end{align*}

Let's consider the term $\langle V(\tau')V(\tau) \rangle - \langle V(\tau') \rangle\langle V(\tau) \rangle$
in the double integral now. This time we have terms with two vertices, and recalling the
linked cluster theorem the expression factorizes as follows:
\begin{align*}
	\langle V(\tau')V(\tau) \rangle - \langle V(\tau') \rangle\langle V(\tau) \rangle
	&= \mathbb{E}_G(V(\tau')V(\tau)) - \mathbb{E}_G(V(\tau'))\mathbb{E}_G(V(\tau)) \\
	&= \mathbb{E}^{\mathrm{T}}_G(V(\tau');V(\tau))\\
	&= \sum_{\Gamma\in\mathcal{G}_{\mathrm{C},2}} \mathrm{val}(\Gamma) = \star
\end{align*}

Assuming that $V$ is such that $\hat{V}(-k)=\hat{V}(k)$, the
connected diagrams with two vertices are exactly divided in two groups, up to renaming the
indices in the resulting sums. Indeed, on one hand we have diagrams that have no loops, i.e.
each vertex only has contractions with the other vertex and never with itself, and on the
other we have diagrams with two loops.

\textcolor{red}{draw the complete set of diagrams?}.

Consider the first group. Fixing an incoming and an outgoing leg on one vertex, there are 4 ways
to contract them with legs of the other vertex, and in each case there is only one way to
complete the diagram without introducing loops, therefore this group contains 4 diagrams.
For the second group, there are 4 ways to form a loop on a vertex, and given a (connected)
diagram with two vertices and two loops there exists only one way to complete it. Therefore
the second group contains 16 diagrams (all possible ways to put one loop on both vertices).
Accounting for these multiplicities, we have:
\begin{align*}
	\star &= 4 \left( \textcolor{red}{no loops} \right) + 16 \left( \textcolor{red}{with loops} \right)\\
	&= \sum_{k,p,q,k',p',q'} \hat{V}(k')\hat{V}(k) \left(
		4 \, G_{p-k,p'}G_{q'+k',q}G_{q+k,q'}G_{p'-k',p}
		+ 16 \, G_{p-k,q}G_{q+k,q'}G_{p'-k',p}G_{q'+k',p'}
	\right) \\
	&= \sum_{p,q,p'} \Bigl(
		4 \, (\hat{V}(p-p'))^2 e^{(\tau'-\tau)(\epsilon_p + \epsilon_q + \epsilon_{p'} + \epsilon_{p-p'+q})} n_p n_q n_{p'} n_{p-p'+q} \\
	&\qquad\qquad + 16 \, \hat{V}(0)\hat{V}(p-q) e^{(\tau'-\tau)(2\epsilon_p + \epsilon_q + \epsilon_{p'})} n_p^2 n_q n_{p'}
	\Bigr)
\end{align*}

Renaming some of the indices at both orders and integrating this last expression, as per
equation (\ref{eq:duhamel}), leads to the final expression of the free energy at second order

\begin{align*}
	\mathrm{log}\,\mathrm{Tr}\,e^{-\beta H}
	&= \mathrm{log}\,\mathrm{Tr}\,e^{-\beta H_0}\\
	&\quad-\beta \sum_{p,k\in\mathcal{D}_\Lambda} \left( \hat{V}(0) -\hat{V}(k) \right) n_p n_{p-k} \\
	&\quad - \sum_{p,q,k\in\mathcal{D}_\Lambda} \Biggr(
	\frac{4(\hat{V}(k))^2 n_p n_q n_{q+k} n_{p-k}}{\epsilon_p + \epsilon_q + \epsilon_{q+k} + \epsilon_{p-k}} \biggl( \beta + \frac{1 - e^{-\beta (\epsilon_p + \epsilon_q + \epsilon_{q+k} + \epsilon_{p-k})}}{\epsilon_p + \epsilon_q + \epsilon_{q+k} + \epsilon_{p-k}} \biggr) \\
	&\quad +
	\frac{16\hat{V}(0)\hat{V}(p-q) n^2_p n_q n_{p-k}}{2\epsilon_p + \epsilon_q + \epsilon_{p-k}} \biggl( \beta + \frac{1 - e^{-\beta (2\epsilon_p + \epsilon_q + \epsilon_{p-k})}}{2\epsilon_p + \epsilon_q + \epsilon_{p-k}} \biggr)
	\Biggl).
\end{align*}

Now that the computation is complete, we need to prove that removing the UV cutoff, i.e.
taking the limit $\Lambda\to\infty$, leads to a well-defined expression. Please note that this
is not at all a trivial statement and what we are doing is deciding on the order in which
to take different limits, as removing the cutoff starting from this expression means that the
perturbative series is \textit{not} done in the case where the state space is
infinite-dimensional. Physically, this does not pose any concern and we expect the two
scenarios to be equivalent, but from the point of view of mathematical rigour it is worth
pointing out this kind of subtleties.

\begin{theorem}
	With the regularity condition that, given $m\geq1$, $\partial^\alpha V\in L^1(\mathbb{T}^3)$
	for all multi-indices $\alpha$ such that $\vert\alpha\vert=m$, and the assumption that
	$\epsilon_p$ is a positive and monotonically increasing \textcolor{red}{polynomial} in
	$\vert p\vert$, the above expression converges in the limit $\Lambda\to\infty$.
\end{theorem}
\begin{proof}
	The zero-th order is trivial, as it doesn't depend on $\Lambda$ and is, in fact, finite:
	$e^{-\beta H_0}$ is trace-class and strictly positive definite, which implies
	$\vert\log\,\mathrm{Tr}\,e^{-\beta H_0}\vert<\infty$.
	Moving onto the first and second orders, by the non-stationary phase argument we can
	conclude that $\vert\hat{V}(k)\vert \leq C/\vert k\vert^m$ for some $C>0$, which means
	that at both orders the factors $\hat{V}(k)$ will be decaying as an inverse power as
	$\vert k\vert\to\infty$, which is the relevant limit as we take the cutoff $\Lambda$ to
	infinity. As for the energy terms, we are assuming that $\epsilon_k\geq0$ for all $k$, and
	that it is a monotonically increasing polynomial in $\vert k\vert$. The last factor to
	discuss is the Fermi-Dirac distribution, but this has an explicit form: an exponential
	decay in $\epsilon_k$, which leads to an explonential decay in some power (greater or
	equal to one) of $\vert k\vert$.

	With these estimates in place we can show that, in the limit $\Lambda\to\infty$, the
	expression above converges. Looking at the first order, we can write
	\begin{align*}
		\sum_{p,k\in\mathcal{D}_\Lambda} \left\vert\left(\hat{V}(0)-\hat{V}(k)\right) n_p n_{p-k}\right\vert
		&\leq \vert\hat{V}(0)\vert \sum_{p,k\in\mathcal{D}_\Lambda}n_p n_{p-k}
		+ \sum_{p,k\in\mathcal{D}_\Lambda}\vert\hat{V}(k)\vert n_p n_{p-k}\\
		&= \vert\hat{V}(0)\vert \sum_{p\in\mathcal{D}_\Lambda}n_p\sum_{k\in\mathcal{D}_\Lambda}n_{p-k}
		+ \sum_{p\in\mathcal{D}_\Lambda}n_p\sum_{k\in\mathcal{D}_\Lambda}\vert\hat{V}(k)\vert n_{p-k}.
	\end{align*}

	By the positivity of the general terms for all $k\in\mathbb{Z}^3$ we can bound the sums by
	the respective sum over $\mathbb{Z}^3$, and because the general terms of each of the above
	sums only depends on the modulus of the index we can employ the following strategy. Define
	$A_n = \{k\in\mathbb{Z}^3 : n \leq \vert k\vert < n+1\}$ and write
	\begin{align*}
		\sum_{k\in\mathbb{Z}^3} a_{\vert k\vert} = \sum_{n=0}^\infty \sum_{k\in A_n} a_{\vert k\vert}.
	\end{align*}
	Because $c_1\vert A_n\vert a_{n+1} \leq \sum_{k\in A_n}a_{\vert k\vert} \leq c_2\vert A_{n+1}\vert a_n$
	for all $n$ if $a_n$ is decreasing, and because $\vert A_n\vert\asymp n^2$, we can conclude
	that the behaviour of $\sum_{k\in\mathbb{Z}^3} a_{\vert k\vert}$ and $\sum_{n=0}^\infty n^2 a_n$
	is the same.

	Now we can show that each term converges. Indeed $\sum_{k\in\mathbb{Z}^3}n_k$ converges
	if, and only if, the series $\sum_{\ell\in\mathbb{N}} \ell^2 e^{-\beta\epsilon_\ell}$
	converges, by the asymptotic comparison test. The latter does converge absolutely as shown by
	the root test:
	\begin{align*}
		\left\vert \ell^2 e^{-\beta\epsilon_{\ell}}\right\vert^{1/\ell}
		= e^{\frac{2\log(\ell)}{\ell} - \beta\epsilon_\ell}
		\longrightarrow K \in [0,1).
	\end{align*}

	Similarly, $\sum_{k\in\mathbb{Z}^3}\vert\hat{V}(k)\vert n_k$ converges, as
	\begin{align*}
		\sum_{k\in\mathbb{Z}^3} \vert\hat{V}(k)\vert n_k
		\leq \sum_{k\in\mathbb{Z}^3}\frac{C n_k}{\vert k\vert^m}
		\leq \sum_{k\in\mathbb{Z}^3} C n_k.
	\end{align*}
	This shows that indeed the first order is, in the limit in $\Lambda\to\infty$, an
	absolutely convergent series.

	Very similar arguments can be used to show that the second order has an absolutely convergent
	series as a limit. Indeed, recalling the assumptions made about each factor present in the
	expression, we can bound the first term in the sum by
	\begin{align*}
		\sum_{p,q,k\in\mathcal{D}_\Lambda} (\ldots)
		&\leq \sum_p \frac{4n_p}{\epsilon_p} \sum_q n_q \sum_k (\hat{V}(k))^2 n_{q+k}n_{p-k}
		\left(\beta + \frac{1}{\epsilon_p}\right),
	\end{align*}
	which, turning once more sums over $\mathbb{Z}^3$ into normal series with an additional square
	in the general term, is an absolutely convergent series in the limit $\Lambda\to\infty$, as
	the only true difference from the previous case is the additional factor $1/\epsilon_p$, which
	is an inverse power and does not change the result of any of the comparison tests done before.
	The exact same argument can be applied to the second term, and this concludes the proof.
\end{proof}

And with the convergence of the free energy under control, the next step is to look at the
limit $\beta\to\infty$ to retrieve the ground state energy.

\subsection{Mean field scaling limit}
